\section{知识共享与未来迭代}

\subsection{知识流与复现}
AlphaProof 团队把 experiment log、failed proof trace、human comment、newsletter 实验记入单一 repo；每个 release 都附带 `diff`，用来 retrain policy。Hubert 指出：“我们需要把 QA 亮点、governance decision、正负例 proof 共享给新成员，这样 training loop 才能迭代得更快。”
\begin{importantbox}{知识共享机制}
1）Experiment repo，包含 chunk gating journal；2）Newsletter，每周记录 QA + governance highlight；3）Shared slides + transcript，也是新成员 ramp-up 的料。
\end{importantbox}

\subsection{问答亮点}
QA 环节被用来说明是否能将 AlphaProof 扩展到 non-Euclidean geometry 以及 ~hypergraph theorem。Hubert 回答时强调要把 RL policy 设计为 modular，每次只调 search depth 或 tactic mix，就可以逐步扩展到更抽象证明。
\begin{warningbox}{问答中的谨慎点}
不要盲目用 AlphaProof 解决所有问题；在 low-resource 语言或 novel math branch 里，policy 需要重新 warm start，并配合 chunk gating 与 fairness board。
\end{warningbox}

\subsection{本章小结}
知识共享段落展示了 log/newsletter/transcript 如何形成复现网络，以及通过 QA 环节暴露的扩展挑战，进一步说明 iteration 的边界与治理需求。
\newpage

